\documentclass[11pt]{amsart}
\usepackage[margin=1in]{geometry}
\usepackage[T1]{fontenc}
\usepackage{lmodern}
\usepackage{amsmath,amssymb,amsthm,mathtools,microtype}
\setlength{\emergencystretch}{2em}
\newtheorem{theorem}{Theorem}
\newcommand{\E}{\mathbb E}
\newcommand{\one}{\mathbf 1}
\pagestyle{plain}

\begin{document}

\begin{theorem}[Two-point reduction]\label{thm:two-point}
Let $\frac12\le\beta\le1$ and let $B:[0,1]\to[0,\infty)$ satisfy
$B(0)=B(1)=0$. Suppose that, for every $0\le x\le y\le1$,
\begin{equation}\label{eq:two-point}
 B(x)+\max\left\{
 \bigl((y-x)^{1/\beta}+B(y)^{1/\beta}\bigr)^\beta,\,
 y-x+(2^\beta-1)B(y)
 \right\}
 \ge 2B\!\left(\frac{x+y}{2}\right).
\end{equation}
Then, for every $n\ge1$ and every $A\subseteq\{0,1\}^n$,
\[
 \E_n h_A^\beta\ge B(|A|).
\]
Here $|A|=2^{-n}\#A$ and $\E_n$ denotes expectation with respect to
normalized counting measure on $\{0,1\}^n$. For $z\in A$, the value
$h_A(z)$ is the number of neighbors of $z$ outside $A$, where two
vertices are neighbors if they differ in exactly one coordinate;
set $h_A(z)=0$ for $z\notin A$.
\end{theorem}

\begin{proof}
Set
\[
 q=\frac1\beta,\qquad c=2^\beta-1,\qquad
 N(s,t)=(s^q+t^q)^{1/q}\quad(s,t\ge0).
\]
Since $q\ge1$, the function $N$ is the $\ell^q$ norm on the
nonnegative quadrant, and $0<c\le1$.

We first establish an estimate for adding an indicator. Let $h$ be
a nonnegative integer-valued function on a finite probability space,
let $D$ be an event, and write $H=\E h^\beta$ and $d=\mathbb P(D)$.
We claim that
\begin{equation}\label{eq:increment}
 \E(h+\one_D)^\beta\ge\max\{N(d,H),\,d+cH\}.
\end{equation}
The triangle inequality for an average of vectors gives
\[
 N(H,d)
 =\bigl\|\E(h^\beta,\one_D)\bigr\|_{\ell^q}
 \le \E\bigl\|(h^\beta,\one_D)\bigr\|_{\ell^q}
 =\E(h+\one_D)^\beta,
\]
which proves the first lower bound.

For the second lower bound, observe that
\begin{equation}\label{eq:scalar}
 (k+1)^\beta\ge1+ck^\beta
 \qquad(k=0,1,2,\ldots).
\end{equation}
This is an equality for $k=0$. If $k\ge1$, set $t=k^\beta\ge1$.
The triangle inequality, applied to $(t,t)=(1,t)+(t-1,0)$, yields
\[
 2^\beta t=\|(t,t)\|_{\ell^q}
 \le\|(1,t)\|_{\ell^q}+t-1
 =(k+1)^\beta+t-1.
\]
Rearranging proves \eqref{eq:scalar}. Applying \eqref{eq:scalar}
on $D$ and using $h^\beta\ge ch^\beta$ on $D^c$, we obtain
\[
 (h+\one_D)^\beta\ge\one_D+ch^\beta.
\]
Taking expectations proves the second lower bound in
\eqref{eq:increment}.

We now prove the theorem by induction on the dimension, allowing
dimension zero for the base case. The zero-dimensional cube consists
of one point and has no edges. Its only subsets have measure zero
or one, so the desired bound follows from $B(0)=B(1)=0$.

Suppose the bound holds in dimension $n$, and let
$A\subseteq\{0,1\}^{n+1}$. Define its slices by
\[
 A_i=\{z\in\{0,1\}^n:(z,i)\in A\},\qquad
 a_i=|A_i|,\qquad H_i=\E_n h_{A_i}^\beta
 \quad(i=0,1).
\]
Relabel the slices so that $a_0\le a_1$. Then
\[
 |A|=\frac{a_0+a_1}{2},
 \qquad H_i\ge B(a_i)
\]
by the induction hypothesis. The boundary degrees satisfy the exact
identity
\[
 h_A(z,i)=h_{A_i}(z)+\one_{A_i\setminus A_{1-i}}(z).
\]
Indeed, the first term counts the outgoing edges in the first $n$
coordinates and the second counts the outgoing edge in the last
coordinate. Both sides vanish when $z\notin A_i$.

The contribution from slice zero is at least $H_0$. Applying
\eqref{eq:increment} to slice one, with $h=h_{A_1}$ and
$D=A_1\setminus A_0$, gives
\[
 2\E_{n+1}h_A^\beta
 \ge H_0+\max\{N(d,H_1),\,d+cH_1\},
 \qquad d=|A_1\setminus A_0|.
\]
Since
\[
 d=a_1-|A_0\cap A_1|\ge a_1-a_0
\]
and both entries in the maximum are nondecreasing in each variable,
the induction hypothesis implies
\begin{align*}
 2\E_{n+1}h_A^\beta
 &\ge B(a_0)+
 \max\{N(a_1-a_0,B(a_1)),\,a_1-a_0+cB(a_1)\}\\
 &\ge 2B\!\left(\frac{a_0+a_1}{2}\right)
 =2B(|A|),
\end{align*}
where the last inequality is \eqref{eq:two-point}. This completes
the induction.
\end{proof}

\end{document}
